\section{Extensions}
\label{sec:extensions}

\subsection{The Cole--Hopf transformation}
\label{subsec:cole-hopf}

This extension reintroduces viscosity and solves the viscous Burgers'
equation \eqref{eq:viscous-burgers} exactly, using the potential-flow
structure noted in Remark~\ref{rem:vorticity-free}.

\begin{theorem}[Cole--Hopf transformation]
\label{thm:cole-hopf}
Let $u$ solve $u_t+u\,u_x=\nu\,u_{xx}$, and define $\varphi>0$ by
\begin{equation}
  \label{eq:cole-hopf-def}
  u = -2\nu\,\frac{\varphi_x}{\varphi}, \qquad \text{i.e.} \qquad \varphi(x,t) = \exp\left(-\frac{1}{2\nu}\int_0^x u(y,t)\, dy\right).
\end{equation}
Then $\varphi$ solves the heat equation $\varphi_t = \nu\, \varphi_{xx}$.
\end{theorem}

\begin{proof}
Write $\psi = \ln\varphi$, so $u=-2\nu\,\psi_x$. Substituting into the
viscous Burgers' equation and dividing by $-2\nu$ gives
\begin{equation}
  \psi_{xt} - 2\nu\,\psi_x\psi_{xx} - \nu\,\psi_{xxx} = 0,
\end{equation}
which is exactly $\partial_x\big[\psi_t - \nu\,\psi_{xx} - \nu\,(\psi_x)^2\big]=0$.
Hence $\psi_t - \nu\,\psi_{xx} - \nu\,(\psi_x)^2 = C(t)$ for some function
$C$ of $t$ alone; since $u=-2\nu\,\psi_x$ is unchanged if $\psi$ is
shifted by a function of $t$, $\varphi$ may be rescaled by
$\exp\!\big(-\!\int_0^t C\big)$ without affecting $u$, so we may take
$C\equiv0$:
\begin{equation}
  \label{eq:psi-equation}
  \psi_t = \nu\,\psi_{xx} + \nu\,(\psi_x)^2.
\end{equation}
Since $\varphi=e^\psi$ gives $\varphi_t=\varphi\,\psi_t$ and
$\varphi_{xx}=\varphi\big(\psi_{xx}+(\psi_x)^2\big)$, equation
\eqref{eq:psi-equation} is equivalent to $\varphi_t=\nu\,\varphi_{xx}$.
\end{proof}

\subsection{Exact solution via the heat kernel}
\label{subsec:cole-hopf-solution}

By Theorem~\ref{thm:cole-hopf}, solving \eqref{eq:viscous-burgers} reduces
to solving the heat equation for $\varphi$ with initial data
\begin{equation}
  \label{eq:phi0}
  \varphi_0(x) = \exp\left(-\frac{1}{2\nu}\int_0^x u_0(y)\, dy\right),
\end{equation}
which is given by convolution with the heat kernel,
\begin{equation}
  \label{eq:phi-heat-kernel}
  \varphi(x,t) = \int_{\mathbb{R}} \frac{1}{\sqrt{4\pi\nu t}}\,
    \exp\left(-\frac{(x-y)^2}{4\nu t}\right) \varphi_0(y)\, dy.
\end{equation}
Substituting \eqref{eq:phi0}--\eqref{eq:phi-heat-kernel} into
$u=-2\nu\,\varphi_x/\varphi$ and differentiating under the integral gives
the explicit \emph{Cole--Hopf formula}
\begin{equation}
  \label{eq:cole-hopf-formula}
  u(x,t) = \frac{\displaystyle\int_{\mathbb{R}} \frac{x-y}{t}\, \exp\left(-\frac{G(x,y,t)}{2\nu}\right) dy}
                {\displaystyle\int_{\mathbb{R}} \exp\left(-\frac{G(x,y,t)}{2\nu}\right) dy},
  \qquad
  G(x,y,t) = \int_0^y u_0(s)\, ds + \frac{(x-y)^2}{2t}.
\end{equation}

\begin{corollary}[Global smoothness]
\label{cor:global-smoothness}
For every bounded measurable $u_0$ and every $\nu>0$, formula
\eqref{eq:cole-hopf-formula} defines the unique bounded solution of
\eqref{eq:viscous-burgers}, and this solution is $C^\infty$ for all
$t>0$.
\end{corollary}

Unlike the inviscid equation (Proposition~\ref{prop:shock-time}),
solutions of \eqref{eq:viscous-burgers} stay smooth for all time, for
any $\nu>0$ however small: advection and viscous diffusion together do
not by themselves produce singularities.

\subsection{The vanishing viscosity limit}
\label{subsec:vanishing-viscosity}

The Cole--Hopf solution also explains why the Lax entropy condition of
Section~\ref{sec:entropy-condition} is the physically correct selection
principle, rather than an arbitrary one.

\begin{theorem}[Vanishing viscosity limit; Hopf, 1950]
\label{thm:vanishing-viscosity}
Let $u_0\in L^\infty(\mathbb{R})$, let $u^\nu$ denote the solution of
\eqref{eq:viscous-burgers} with data $u_0$ given by
\eqref{eq:cole-hopf-formula}, and let $u$ denote the entropy solution of
\eqref{eq:inviscid-burgers} with the same data
(Theorem~\ref{thm:riemann-solution}, extended to general $u_0$ by the same
selection principle). Then $u^\nu \to u$ pointwise almost everywhere,
and in $L^1_{\mathrm{loc}}$, as $\nu \to 0^+$ \cite{hopf1950}.
\end{theorem}

The proof (a Laplace-method asymptotic analysis of
\eqref{eq:cole-hopf-formula} as $\nu\to0$, showing the integral
concentrates near the $y$ minimising $G(x,y,t)$) is not reproduced here.
Its significance is conceptual: the entropy condition of
Definition~\ref{def:lax-entropy} is not an ad hoc rule imposed to force
uniqueness, but exactly the criterion selected by the physically natural
limit of vanishing viscosity --- the same limit by which the inviscid
equation \eqref{eq:inviscid-burgers} was justified as an idealisation of
\eqref{eq:viscous-burgers} in Section~\ref{sec:motivation}.

\subsection{A numerical scheme: Godunov's method}
\label{subsec:godunov}

This extension outlines a numerical scheme for the conservation law
$u_t+f(u)_x=0$ that uses the exact Riemann solution of
Theorem~\ref{thm:riemann-solution} as its building block, and states its
convergence to the entropy solution.

\subsubsection{Finite volume discretisation}
\label{subsubsec:finite-volume}

Divide $\mathbb{R}$ into cells $[x_{j-1/2}, x_{j+1/2}]$ of width $\Delta
x$, with cell averages
\begin{equation}
  U_j^n \approx \frac{1}{\Delta x}\int_{x_{j-1/2}}^{x_{j+1/2}} u(x, t^n)\, dx,
\end{equation}
at times $t^n = n\Delta t$. Integrating $u_t+f(u)_x=0$ over a cell and
over $[t^n,t^{n+1}]$ gives the exact identity
\begin{equation}
  U_j^{n+1} = U_j^n - \frac{\Delta t}{\Delta x}
  \left( \widehat{f}_{j+1/2} - \widehat{f}_{j-1/2} \right),
\end{equation}
where $\widehat f_{j+1/2}$ is the (time-averaged) flux through the
interface $x_{j+1/2}$, which must be approximated by a \emph{numerical
flux function} $\widehat f_{j+1/2} = \widehat f(U_j^n, U_{j+1}^n)$.

\subsubsection{The Godunov flux}
\label{subsubsec:godunov-flux}

Godunov's method defines $\widehat f$ by solving, at each interface, the
Riemann problem with left and right states $U_j^n$ and $U_{j+1}^n$
\emph{exactly}, using Theorem~\ref{thm:riemann-solution}, and evaluating
the resulting self-similar solution along $x/t=0$:
\begin{equation}
  \label{eq:godunov-flux}
  \widehat f_{j+1/2} = f\big(u^R(0;\, U_j^n, U_{j+1}^n)\big),
\end{equation}
where $u^R(\xi; u_\ell, u_r)$ denotes the value at self-similar
coordinate $\xi=x/t$ of the entropy solution from
Theorem~\ref{thm:riemann-solution} with data $(u_\ell,u_r)$. Concretely,
by \eqref{eq:riemann-shock}--\eqref{eq:riemann-rarefaction},
\begin{equation}
  u^R(0; u_\ell,u_r) =
  \begin{cases}
    u_\ell, & u_\ell>u_r,\ s=\tfrac{u_\ell+u_r}{2}>0, \\
    u_r, & u_\ell>u_r,\ s<0, \\
    u_\ell, & u_\ell\le u_r,\ u_\ell\ge0, \\
    u_r, & u_\ell\le u_r,\ u_r\le0, \\
    0, & u_\ell < 0 < u_r,
  \end{cases}
\end{equation}
the last case corresponding to $\xi=0$ lying inside the rarefaction fan of
\eqref{eq:riemann-rarefaction}, i.e. $u^R(0)=(f')^{-1}(0)=0$ for Burgers'
flux.

\subsubsection{Stability and convergence}
\label{subsubsec:godunov-convergence}

\begin{proposition}[CFL condition]
\label{prop:cfl}
The Godunov scheme is stable provided the time step satisfies
\begin{equation}
  \label{eq:cfl}
  \Delta t \le \frac{\Delta x}{\max_j |f'(U_j^n)|}.
\end{equation}
\end{proposition}

Condition \eqref{eq:cfl} (the Courant--Friedrichs--Lewy condition)
ensures that information does not propagate across more than one cell
per time step, consistent with the finite propagation speed of
characteristics established in Section~\ref{sec:characteristics}.

\begin{theorem}[Convergence; see \cite{evans2010} and standard finite
volume references]
\label{thm:godunov-convergence}
Under \eqref{eq:cfl}, the Godunov approximations converge, as $\Delta x,
\Delta t \to 0$, to the unique entropy solution of
Theorem~\ref{thm:oleinik-uniqueness}.
\end{theorem}

\begin{remark}
This gives a second, independent route to the solutions constructed by
hand in Sections~\ref{sec:nonuniqueness}--\ref{sec:riemann-problem}: an
implementation of \eqref{eq:godunov-flux} applied to the Riemann data
\eqref{eq:riemann-data} should reproduce the shock speed $s=1/2$ and the
linear profile $u=x/t$ inside the rarefaction fan, providing a
computational check on the exact formulas derived analytically above.
\end{remark}

\subsubsection{Numerical results}
\label{subsubsec:godunov-results}

The scheme of Section~\ref{subsubsec:finite-volume}--\ref{subsubsec:godunov-flux}
was implemented with $N=400$ cells on $x\in[-2,2]$, using the CFL number
$0.45$ in \eqref{eq:cfl}, and run to $t=1$ on both Riemann problems
considered earlier: the rarefaction data $(u_\ell,u_r)=(0,1)$ of
Section~\ref{sec:nonuniqueness}, and the shock data $(u_\ell,u_r)=(1,0)$.
Figure~\ref{fig:godunov-results} compares the numerical solution against
the exact entropy solution of Theorem~\ref{thm:riemann-solution}.

\begin{figure}[h]
\centering
\includegraphics[width=0.44\textwidth]{figures/godunov_rarefaction.pdf}
\includegraphics[width=0.44\textwidth]{figures/godunov_shock.pdf}
\caption{Godunov scheme (blue dots, $N=400$ cells) against the exact
entropy solution (red) of Theorem~\ref{thm:riemann-solution}, at
$t=1$. Left: rarefaction data $(u_\ell,u_r)=(0,1)$; the scheme
reproduces the linear fan $u=x/t$, with mild smearing at the two corners
typical of a first-order scheme. Right: shock data $(u_\ell,u_r)=(1,0)$;
the discontinuity is captured sharply at $x=st=0.5$, exactly the
Rankine--Hugoniot speed from \eqref{eq:rankine-hugoniot-burgers}, with
no spurious oscillation.}
\label{fig:godunov-results}
\end{figure}

The measured $L^1$ errors between the numerical and exact solutions are
approximately $4.5\times10^{-3}$ (rarefaction case) and $1.2\times10^{-3}$
(shock case), consistent with the first-order convergence rate of
Theorem~\ref{thm:godunov-convergence}; the shock case converges faster
per cell because the error is concentrated in the single cell straddling
the discontinuity, whereas the rarefaction error accumulates gradually
across the two fan corners.

\subsection{Long-time behaviour: the N-wave}
\label{subsec:n-wave}

This final extension describes the universal shape that entropy
solutions with compactly supported initial data approach as $t \to
\infty$, regardless of the detailed shape of $u_0$.

\subsubsection{An exact self-similar solution}
\label{subsubsec:n-wave-exact}

\begin{proposition}[The N-wave]
\label{prop:n-wave-exact}
Fix constants $P, Q > 0$. The function
\begin{equation}
  \label{eq:n-wave}
  N(x,t) = \begin{cases}
    0, & x \le -\sqrt{2Qt}, \\
    x/t, & -\sqrt{2Qt} < x < \sqrt{2Pt}, \\
    0, & x \ge \sqrt{2Pt},
  \end{cases}
  \qquad t>0,
\end{equation}
is an entropy solution of $u_t+u\,u_x=0$ on $\mathbb{R}\times(0,\infty)$.
\end{proposition}

\begin{proof}
On the open middle interval, $N(x,t)=x/t$ is exactly the self-similar
rarefaction profile of Section~\ref{subsec:rarefaction-general}, which
solves the PDE classically there; the outer regions are trivially
constant solutions. At the right edge $x_+(t)=\sqrt{2Pt}$, the one-sided
limits are $u_\ell = x_+(t)/t = \sqrt{2P/t}$ and $u_r=0$; since
$\dot x_+(t) = \sqrt{P/(2t)} = \tfrac12\big(\sqrt{2P/t}+0\big)$, this is
exactly the Rankine--Hugoniot speed \eqref{eq:rankine-hugoniot-burgers},
and $u_\ell > u_r$ so the Lax condition \eqref{eq:lax-burgers} holds. The
same check at the left edge $x_-(t)=-\sqrt{2Qt}$, where $u_\ell=0$ and
$u_r=-\sqrt{2Q/t}$, gives $\dot x_-(t) = -\sqrt{Q/(2t)} =
\tfrac12\big(0+(-\sqrt{2Q/t})\big)$, again matching
\eqref{eq:rankine-hugoniot-burgers}, with $u_\ell=0>-\sqrt{2Q/t}=u_r$
satisfying the Lax condition. Hence both jumps are entropy-admissible,
and $N$ is an entropy solution by Theorem~\ref{thm:rankine-hugoniot} and
Definition~\ref{def:lax-entropy}.
\end{proof}

The profile \eqref{eq:n-wave} is called an \emph{N-wave}: it rises
linearly through the origin and is bounded by a shock on each side, so
that its graph traces the shape of the letter N (reversed), decaying in
amplitude like $t^{-1/2}$ and spreading in width like $t^{1/2}$ as
$t\to\infty$. It is an exact solution for every $t>0$, though its formal
limit as $t\to0^+$ is not an ordinary function of the type considered in
Section~\ref{sec:weak-solutions}: both edges collapse to the origin,
consistent with $N(\cdot,t)$ representing, in the limit, a state
concentrated entirely at a point.

\subsubsection{Shock merging}
\label{subsubsec:shock-merging}

The convergence result below relies on one further fact: colliding
shocks merge into one, governed by exactly the Riemann-problem machinery
already developed.

\begin{lemma}[Shock merging]
\label{lem:shock-merging}
Let $u_1>u_2>u_3$, and consider initial data equal to $u_1$ for $x<a$,
$u_2$ for $a<x<b$, and $u_3$ for $x>b$. For small $t>0$ the entropy
solution consists of two shocks moving at speeds $s_{12}=(u_1+u_2)/2$
and $s_{23}=(u_2+u_3)/2$; since $s_{12}>s_{23}$, they collide at the
finite time $t^\times = 2(b-a)/(u_1-u_3)$, after which the entropy
solution is the single shock connecting $u_1$ to $u_3$ at speed
$s_{13}=(u_1+u_3)/2$.
\end{lemma}

\begin{proof}[Proof sketch]
Each jump individually satisfies the Lax condition since
$u_1>u_2>u_3$, so Theorem~\ref{thm:riemann-solution}(a) gives the two
shock speeds above. Since $s_{12}-s_{23}=\tfrac12(u_1-u_3)>0$, the shock
positions $a+s_{12}t$ and $b+s_{23}t$ meet at the stated finite
$t^\times$. At that instant the middle state has shrunk to zero width,
leaving the two-state data $u_1,u_3$; since $u_1>u_3$,
Theorem~\ref{thm:riemann-solution}(a) applied afresh at $(x^\times,
t^\times)$ gives the merged shock at speed $s_{13}$.
\end{proof}

\subsubsection{Convergence of general solutions to the N-wave}
\label{subsubsec:n-wave-convergence}

\begin{theorem}[Asymptotic N-wave decay; Lax, 1957 \cite{lax1957}]
\label{thm:n-wave-decay}
Let $u_0 \in L^1(\mathbb{R})$ have compact support, and set
\begin{equation}
  \label{eq:pq-areas}
  P = \int_{\mathbb{R}} \max(u_0(x),0)\, dx, \qquad
  Q = \int_{\mathbb{R}} \max(-u_0(x),0)\, dx.
\end{equation}
Then the entropy solution $u$ of \eqref{eq:ivp} satisfies
\begin{equation}
  \|u(\cdot,t) - N(\cdot,t)\|_{L^1(\mathbb{R})} \to 0 \qquad \text{as } t \to \infty,
\end{equation}
where $N$ is the N-wave \eqref{eq:n-wave} with the same constants $P,Q$.
\end{theorem}

The proof is not reproduced here, but the mechanism is exactly
Lemma~\ref{lem:shock-merging}: a compactly supported $u_0$ produces, for
small $t$, a solution with possibly many shocks and rarefactions, but
every pair of shocks that meet coalesces into one, so only finitely many
discontinuities survive as $t\to\infty$. Since the total signed area
$\int u(x,t)\,dx = P-Q$ is conserved for all $t$ (integrate
\eqref{eq:weak-formulation} over all $x$), the surviving structure is
forced into the unique self-similar shape consistent with these boundary
areas --- exactly \eqref{eq:n-wave}.